% Part VII -- Algorithms, Optimization, and Machine Learning
% Chapter 22: Neural networks

\h{Neural Networks}
\label{ch:22}

A neural network is not a new idea bolted onto the previous chapters. It is the
two ideas you already have, combined: a model whose parameters are chosen by
gradient descent (\booklink{ch:20}{Chapter 20}) and whose architecture is
selected by cross-validation (\booklink{ch:21}{Chapter 21}). What is new is that
the model composes several layers of weighted sums, and that a nonlinear
function between those layers is what makes the composition worth doing. This
chapter builds one network from arrays, checks its gradient, trains it, and then
asks the question that matters in practice: on this problem, is it better than
what you already had?

\begin{NoteBox}
\textbf{Prerequisites.} You need NumPy array shapes and matrix multiplication,
gradient descent and finite-difference gradient checking from
\booklink{ch:20}{Chapter 20}, and the classification workflow of
\booklink{ch:21}{Chapter 21}: stratified splits, pipelines, cross-validation,
ROC-AUC, and validation curves. This chapter reuses that chapter's
comfort-complaint dataset so the comparison is direct.
\end{NoteBox}

\begin{tcolorbox}[title=Learning outcomes,colframe=orange,colback=white,arc=2mm]
After completing this chapter, you should be able to:
\begin{itemize}
    \item describe a neuron as a weighted sum followed by an activation;
    \item demonstrate that stacked linear layers collapse into one linear layer;
    \item compare sigmoid, tanh, and ReLU by their output and their slope;
    \item write a forward pass with matrix multiplication and verify its shapes;
    \item explain why binary cross-entropy, not accuracy, is the training loss;
    \item derive and implement backpropagation for a one-hidden-layer network;
    \item verify backpropagation against a finite-difference gradient;
    \item train a network and read its loss curve;
    \item explain why feature scaling is mandatory rather than advisable;
    \item tune \c{MLPClassifier} with validation and learning curves; and
    \item judge honestly when a neural network is not the right model.
\end{itemize}
\end{tcolorbox}

\hh{A neuron is a weighted sum and an activation}

One neuron takes several inputs, multiplies each by a learned weight, adds a
learned bias, and passes the total through an activation function.

\bookfigure[0.98\textwidth]{Figures/Generated/ch22_neuron_anatomy.pdf}
{A neuron multiplies each input by a learned weight, adds a bias, and applies an activation; without that activation the entire network reduces to a single weighted sum.}
{fig:ch22-neuron}

\[
a = g\!\left(w_1x_1 + w_2x_2 + \cdots + w_px_p + b\right).
\]

If $g$ is the sigmoid function and there is only one such neuron, this is exactly
logistic regression from \booklink{ch:21}{Chapter 21}. A network is what you get
by putting many neurons side by side into a \textbf{layer}, and stacking layers
so one layer's outputs are the next layer's inputs.

\hh{Without a nonlinearity, depth buys nothing}

The claim that activations are essential is not a matter of taste. It is
arithmetic, and you can check it in three lines.

\begin{lstlisting}[
    language=python,
    style=block,
    caption={Two stacked linear layers are one linear layer},
    label={c22_collapse},
    captionpos=b
]
import numpy as np

rng = np.random.default_rng(3)
first_layer = rng.normal(size=(6, 4))
second_layer = rng.normal(size=(4, 1))
sample = rng.normal(size=(5, 6))

two_layers = sample @ first_layer @ second_layer
one_layer = sample @ (first_layer @ second_layer)

print("Identical:", np.allclose(two_layers, one_layer))
print("Combined weight shape:", (first_layer @ second_layer).shape)
\end{lstlisting}

Expected output:

\begin{lstlisting}[language=text, style=block, caption={Depth without nonlinearity is an illusion}, label={c22_collapse_output}, captionpos=b]
Identical: True
Combined weight shape: (6, 1)
\end{lstlisting}

\begin{ExplainBox}{Why the activation is the whole point}
Matrix multiplication is associative, so a chain of weight matrices can always
be multiplied together in advance into one matrix of the same overall shape. A
hundred stacked linear layers therefore express precisely the set of functions
one linear layer expresses -- with a hundred times the parameters and none of the
extra capability. Inserting a nonlinear $g$ between layers breaks the
associativity, and the composition can then bend, fold, and enclose regions of
the input space that no straight boundary can separate.
\end{ExplainBox}

\hh{Activation functions differ in their slope, not only their shape}

\bookfigure[0.98\textwidth]{Figures/Generated/ch22_activations.pdf}
{Sigmoid and tanh flatten at both extremes so their slope approaches zero and learning stalls in the shaded regions; ReLU keeps a slope of one for positive inputs but a slope of exactly zero below.}
{fig:ch22-activations}

\begin{center}
\begin{tabular}{lll}
\textbf{Activation} & \textbf{Output range} & \textbf{Practical note} \\
Sigmoid $\sigma(z)$ & $(0, 1)$ & Maximum slope 0.25; use it for the output of a binary classifier \\
Tanh $\tanh(z)$ & $(-1, 1)$ & Zero-centred, maximum slope 1; a reasonable hidden default \\
ReLU $\max(0, z)$ & $[0, \infty)$ & Slope 1 or 0; standard in deep networks, and cheap \\
\end{tabular}
\end{center}

The right panel of Figure~\ref{fig:ch22-activations} matters more than the left.
Training multiplies these slopes together layer by layer, so an activation whose
slope is near zero passes almost no signal back to the earlier layers. That is
the \textbf{vanishing gradient} problem, and it is why sigmoid, whose slope never
exceeds 0.25, fell out of favour for hidden layers. This chapter uses tanh in the
hidden layer and sigmoid only at the output, where a probability is required.

\hh{A layer is a matrix multiplication, so watch the shapes}

Writing the layer as a matrix product lets NumPy process every sample at once.

\begin{center}
\begin{tabular}{lll}
\textbf{Object} & \textbf{Shape} & \textbf{Meaning} \\
\c{X} & (n, 6) & n samples, 6 features \\
\c{W1} & (6, 8) & one column of weights per hidden neuron \\
\c{b1} & (8,) & one bias per hidden neuron \\
\c{A1} & (n, 8) & hidden output for every sample \\
\c{W2} & (8, 1) & one weight per hidden neuron into the output \\
\c{probabilities} & (n,) & one probability per sample \\
\end{tabular}
\end{center}

Almost every beginner error in a hand-written network is a shape error. Print the
shapes before printing the values.

\hh{Prepare the Chapter 21 data for a network}

The dataset is the comfort-complaint generator from
Listing~\ref{c21_data}, unchanged, so every result in this chapter can be
compared directly with the previous one. It is repeated here so this chapter
runs on its own.

\begin{lstlisting}[
    language=python,
    style=block,
    caption={Split, then scale using training statistics only},
    label={c22_data},
    captionpos=b
]
import numpy as np
import pandas as pd
from sklearn.model_selection import train_test_split
from sklearn.preprocessing import StandardScaler


def make_comfort_data(seed=42, sample_count=900):
    rng = np.random.default_rng(seed)
    features = pd.DataFrame({
        "floor_area_m2": rng.uniform(60, 520, sample_count),
        "glazing_ratio": rng.uniform(0.10, 0.68, sample_count),
        "occupant_density": rng.uniform(0.02, 0.16, sample_count),
        "setpoint_deviation_c": rng.normal(0.0, 2.4, sample_count),
        "outdoor_temp_c": rng.uniform(-8, 34, sample_count),
        "months_since_service": rng.integers(0, 36, sample_count),
    })
    risk = (
        -9.0
        + 1.10 * np.abs(features["setpoint_deviation_c"])
        + 4.5 * features["glazing_ratio"] * (features["outdoor_temp_c"] > 26)
        + 30.0 * features["occupant_density"]
        + 0.11 * features["months_since_service"]
        + rng.normal(0, 0.30, sample_count)
    )
    probability = 1 / (1 + np.exp(-risk))
    target = pd.Series(
        (rng.random(sample_count) < probability).astype(int),
        name="complaint",
    )
    return features, target


X, y = make_comfort_data()
X_train, X_test, y_train, y_test = train_test_split(
    X, y, test_size=0.25, random_state=42, stratify=y
)

scaler = StandardScaler().fit(X_train)
X_scaled = scaler.transform(X_train)
X_scaled_test = scaler.transform(X_test)
targets = y_train.to_numpy().astype(float)
targets_test = y_test.to_numpy().astype(float)

print("Scaled training shape:", X_scaled.shape)
print("Training column means:", X_scaled.mean(axis=0).round(3))
\end{lstlisting}

Expected output:

\begin{lstlisting}[language=text, style=block, caption={Every feature now has mean zero and unit variance}, label={c22_data_output}, captionpos=b]
Scaled training shape: (675, 6)
Training column means: [ 0. -0. -0.  0.  0.  0.]
\end{lstlisting}

\begin{SyntaxBox}{Fit the scaler on training rows, transform both}
\c{fit()} learns each column's mean and standard deviation, and it sees only
\c{X_train}. \c{transform()} then applies those same training statistics to both
parts, so the test rows are scaled by numbers that were never derived from them.
The targets are converted to a float array because the network's arithmetic mixes
them with predicted probabilities. Later sections replace this manual sequence
with a \c{Pipeline}, which enforces the same order automatically inside every
cross-validation fold.
\end{SyntaxBox}

\hh{The forward pass}

\begin{lstlisting}[
    language=python,
    style=block,
    caption={Initialize a network and compute its predictions},
    label={c22_forward},
    captionpos=b
]
import numpy as np


def sigmoid(values):
    return 1 / (1 + np.exp(-values))


def initialize(input_count, hidden_count, seed=0):
    """Small random weights, zero biases, one reproducible stream."""
    rng = np.random.default_rng(seed)
    return {
        "W1": rng.normal(0, np.sqrt(1 / input_count), (input_count, hidden_count)),
        "b1": np.zeros(hidden_count),
        "W2": rng.normal(0, np.sqrt(1 / hidden_count), (hidden_count, 1)),
        "b2": np.zeros(1),
    }


def forward(parameters, X):
    """Return the hidden activations and the output probabilities."""
    hidden_input = X @ parameters["W1"] + parameters["b1"]
    hidden_output = np.tanh(hidden_input)
    output_input = hidden_output @ parameters["W2"] + parameters["b2"]
    return hidden_output, sigmoid(output_input).ravel()
\end{lstlisting}

\begin{SyntaxBox}{Read the forward pass line by line}
\begin{itemize}
    \item \c{X @ parameters["W1"]} produces an (n, 8) array; adding \c{b1} of
          shape (8,) broadcasts the bias across every row.
    \item \c{np.tanh()} applies elementwise, so the shape is unchanged.
    \item The second product reduces (n, 8) to (n, 1), and \c{ravel()} flattens
          it to (n,) so it aligns with a one-dimensional target.
    \item Weights are initialized randomly but biases at zero. Initializing all
          weights to zero would make every hidden neuron compute the same thing
          and receive the same gradient forever -- they would never differentiate.
    \item The initial standard deviation $\sqrt{1/\text{fan-in}}$ keeps the
          weighted sums small enough that tanh does not start saturated.
\end{itemize}
\end{SyntaxBox}

\hh{Cross-entropy turns predictions into one differentiable number}

Training needs a scalar loss whose gradient exists. Accuracy fails that test: it
counts labels, so it is flat almost everywhere and jumps in steps, and a flat
loss gives gradient descent nothing to descend.

Binary cross-entropy uses the predicted probability directly:

\[
L = -\frac{1}{n}\sum_{i=1}^{n}\Big[y_i\log \hat{y}_i + (1-y_i)\log(1-\hat{y}_i)\Big].
\]

\begin{lstlisting}[
    language=python,
    style=block,
    caption={Binary cross-entropy with a numerical guard},
    label={c22_loss},
    captionpos=b
]
def cross_entropy(probabilities, targets):
    safe = np.clip(probabilities, 1e-12, 1 - 1e-12)
    return -np.mean(
        targets * np.log(safe) + (1 - targets) * np.log(1 - safe)
    )
\end{lstlisting}

When the target is 1 only the first term survives, and it grows without bound as
the predicted probability approaches 0. A confident wrong answer is therefore
punished far more than an uncertain one. The \c{np.clip()} guard exists because
\c{np.log(0.0)} is negative infinity; without it one saturated prediction turns
the whole loss into \c{nan}.

\hh{Backpropagation is the chain rule run backward}

The loss depends on \c{W1} only through a chain: \c{W1} sets the hidden input,
which sets the hidden output, which sets the output input, which sets the
prediction, which sets the loss. The chain rule multiplies the local derivatives
along that path, and backpropagation is the bookkeeping that computes them once
each, from the output backward.

For this network the algebra collapses neatly. Writing $\delta_2$ for the
derivative of the loss with respect to the output's weighted sum:

\begin{align*}
\delta_2 &= \frac{\hat{y} - y}{n}, \\
\frac{\partial L}{\partial W_2} &= A_1^{\top}\delta_2, \qquad
\frac{\partial L}{\partial b_2} = \sum_i \delta_{2,i}, \\
\delta_1 &= \left(\delta_2 W_2^{\top}\right) \odot \left(1 - A_1^{2}\right), \\
\frac{\partial L}{\partial W_1} &= X^{\top}\delta_1, \qquad
\frac{\partial L}{\partial b_1} = \sum_i \delta_{1,i}.
\end{align*}

\begin{lstlisting}[
    language=python,
    style=block,
    caption={Backpropagation for one hidden layer},
    label={c22_backward},
    captionpos=b
]
def backward(parameters, X, targets, hidden_output, probabilities):
    """Return one gradient array per parameter."""
    sample_count = X.shape[0]

    output_delta = ((probabilities - targets) / sample_count).reshape(-1, 1)
    grad_W2 = hidden_output.T @ output_delta
    grad_b2 = output_delta.sum(axis=0)

    hidden_delta = (output_delta @ parameters["W2"].T) * (1 - hidden_output ** 2)
    grad_W1 = X.T @ hidden_delta
    grad_b1 = hidden_delta.sum(axis=0)

    return {"W1": grad_W1, "b1": grad_b1, "W2": grad_W2, "b2": grad_b2}
\end{lstlisting}

\begin{SyntaxBox}{Where each expression comes from}
\begin{itemize}
    \item \c{probabilities - targets} is the whole derivative of cross-entropy
          through the sigmoid. The sigmoid slope
          $\hat{y}(1-\hat{y})$ cancels against the loss derivative, which is why
          this pairing is standard.
    \item Dividing by \c{sample_count} matches the mean in the loss, so the
          gradient scale does not depend on how many rows you trained on.
    \item \c{(1 - hidden_output ** 2)} is the tanh derivative expressed in terms
          of its own output, so no extra array is needed.
    \item Each \c{grad_W} is a transposed input matrix times the delta arriving
          at that layer, and each \c{grad_b} is a column sum of the same delta.
    \item Check the shapes: \c{grad_W1} must match \c{W1} at (6, 8), and
          \c{grad_W2} must match \c{W2} at (8, 1).
\end{itemize}
\end{SyntaxBox}

\hh{Check the gradient before trusting the training loop}

A wrong backpropagation implementation usually still runs, and the loss usually
still decreases a little. \booklink{ch:20}{Chapter 20} gave the remedy: compare
every analytic derivative against a central finite difference.

\begin{lstlisting}[
    language=python,
    style=block,
    caption={Verify backpropagation parameter by parameter},
    label={c22_gradcheck},
    captionpos=b
]
def gradient_check(parameters, X, targets, step=1e-6):
    hidden_output, probabilities = forward(parameters, X)
    analytic = backward(parameters, X, targets, hidden_output, probabilities)

    worst_gap = 0.0
    for key in parameters:
        flat = parameters[key].ravel()
        for index in range(flat.size):
            original = flat[index]
            flat[index] = original + step
            raised = cross_entropy(forward(parameters, X)[1], targets)
            flat[index] = original - step
            lowered = cross_entropy(forward(parameters, X)[1], targets)
            flat[index] = original

            estimate = (raised - lowered) / (2 * step)
            gap = abs(estimate - analytic[key].ravel()[index])
            worst_gap = max(worst_gap, gap)

    return worst_gap


untrained = initialize(X_scaled.shape[1], 8, seed=0)
print(f"Largest gap: {gradient_check(untrained, X_scaled, targets):.3e}")
\end{lstlisting}

On the comfort data with 6 inputs and 8 hidden neurons this prints a largest gap
near \c{1.5e-10}. Anything below roughly \c{1e-7} indicates agreement; a gap near
\c{1e-2} means a real error in the derivation. Note that \c{flat} is a view into
the parameter array, so assigning to \c{flat[index]} modifies the network in
place -- which is why the original value is always restored.

\begin{TryBox}{Break it deliberately}
Remove the \c{(1 - hidden_output ** 2)} factor from \c{hidden_delta} and rerun
the check. The gap should jump by many orders of magnitude. Then train with the
broken gradient and observe that the loss still falls for a while. This is why
the check exists.
\end{TryBox}

\hh{Train the network}

The training loop is gradient descent, unchanged from
\booklink{ch:20}{Chapter 20}: compute the gradient, take a step opposite it,
repeat. One pass over all training rows is called an \textbf{epoch}.

\begin{lstlisting}[
    language=python,
    style=block,
    caption={Full-batch training with a recorded loss history},
    label={c22_train},
    captionpos=b
]
def train(X, targets, hidden_count=8, learning_rate=0.5, epochs=400, seed=0):
    parameters = initialize(X.shape[1], hidden_count, seed)
    history = []

    for _ in range(epochs):
        hidden_output, probabilities = forward(parameters, X)
        gradients = backward(parameters, X, targets, hidden_output, probabilities)
        for key in parameters:
            parameters[key] -= learning_rate * gradients[key]
        history.append(cross_entropy(probabilities, targets))

    return parameters, np.array(history)


parameters, history = train(X_scaled, targets)
print(f"Loss: {history[0]:.4f} -> {history[-1]:.4f}")
\end{lstlisting}

Expected output:

\begin{lstlisting}[language=text, style=block, caption={400 epochs of full-batch descent}, label={c22_train_output}, captionpos=b]
Loss: 0.7663 -> 0.3554
\end{lstlisting}

\bookfigure[0.98\textwidth]{Figures/Generated/ch22_network_training.pdf}
{Cross-entropy falls from 0.766 to 0.355 over 400 epochs with the held-out loss tracking it closely; on two features the network's 0.50 boundary encloses the comfortable middle from both sides while the logistic boundary can only cut the plane with one straight line.}
{fig:ch22-training}

The right panel of Figure~\ref{fig:ch22-training} is the payoff. Complaints rise
when the setpoint deviation is large in \emph{either} direction, so the region of
comfortable rooms is a band in the middle. The network's boundary wraps around
that band from both sides. Logistic regression, restricted to one straight line,
can only shave off a corner -- which is exactly why it scored ROC-AUC 0.776 in
\booklink{ch:21}{Chapter 21} while this network reaches 0.903 on the same test
rows.

\begin{ExplainBox}{Reading a loss curve}
A healthy curve falls steeply and then flattens. A curve that is flat from the
start usually means the learning rate is too small, the features are unscaled, or
the activations are saturated. A curve that oscillates or rises means the
learning rate is too large, exactly as in \booklink{ch:20}{Chapter 20}. A
training curve that keeps falling while the held-out curve turns upward is
overfitting, and the epoch where they separate is where training should have
stopped.
\end{ExplainBox}

\hh{Scaling is not advisable, it is mandatory}

In \booklink{ch:21}{Chapter 21} scaling improved logistic regression's
conditioning. For a network it decides whether training works at all. Compare the
same estimator with and without a \c{StandardScaler} in front of it:

\begin{center}
\begin{tabular}{lll}
\textbf{Estimator} & \textbf{Cross-validated ROC-AUC} & \textbf{Test ROC-AUC} \\
\c{MLPClassifier} alone & 0.617 +/- 0.059 & 0.681 \\
\c{StandardScaler} then \c{MLPClassifier} & 0.837 +/- 0.042 & 0.895 \\
\end{tabular}
\end{center}

\begin{ExplainBox}{Why unscaled features cripple a network}
\c{floor_area_m2} spans hundreds while \c{glazing_ratio} spans less than one. The
initial weighted sums are therefore dominated by the large-scale column and land
far out on the tanh curve, where the slope is near zero and almost no gradient
returns. Meanwhile a single learning rate must serve parameters whose useful step
sizes differ by two orders of magnitude. Standardizing puts every feature on a
comparable scale so one learning rate is appropriate for all of them, and it must
happen inside the \c{Pipeline} so the scaler is refitted within each
cross-validation fold.
\end{ExplainBox}

\hh{Use MLPClassifier for applied work}

Your hand-written network taught you what happens inside. For real use,
scikit-learn's implementation adds mini-batching, adaptive optimizers, weight
regularization, and early stopping.

\begin{lstlisting}[
    language=python,
    style=block,
    caption={A network as one more estimator in the Chapter 21 workflow},
    label={c22_mlp},
    captionpos=b
]
from sklearn.neural_network import MLPClassifier
from sklearn.pipeline import Pipeline
from sklearn.preprocessing import StandardScaler

network = Pipeline([
    ("scale", StandardScaler()),
    ("model", MLPClassifier(
        hidden_layer_sizes=(16,),
        alpha=1e-3,
        max_iter=2000,
        random_state=42,
    )),
])

network.fit(X_train, y_train)
fitted = network.named_steps["model"]

print("Iterations run:", fitted.n_iter_)
print("Final training loss:", round(fitted.loss_, 4))
print("Weight matrix shapes:", [layer.shape for layer in fitted.coefs_])
\end{lstlisting}

Expected output:

\begin{lstlisting}[language=text, style=block, caption={A 129-parameter network}, label={c22_mlp_output}, captionpos=b]
Iterations run: 546
Final training loss: 0.3273
Weight matrix shapes: [(6, 16), (16, 1)]
\end{lstlisting}

\begin{SyntaxBox}{The settings that actually matter}
\begin{itemize}
    \item \c{hidden_layer_sizes=(16,)} is one hidden layer of 16 neurons.
          \c{(16, 8)} would be two layers. It is a tuple, and a common error is
          writing \c{16} instead of \c{(16,)}.
    \item \c{alpha} is L2 weight regularization: larger values shrink weights and
          reduce overfitting. It is a hyperparameter, not a fixed constant.
    \item \c{max_iter} caps the optimizer's iterations. A
          \c{ConvergenceWarning} means it stopped early, and the fix is more
          iterations or better scaling -- never suppressing the warning.
    \item \c{random_state} fixes the weight initialization and batch shuffling.
          Without it, two runs give different models.
    \item \c{early_stopping=True} holds out part of the training data and stops
          when validation stops improving, which is often better than tuning
          \c{max_iter} by hand.
\end{itemize}
Counting parameters: $6 \times 16 + 16$ for the first layer and
$16 \times 1 + 1$ for the second gives 129 numbers fitted from 675 training rows.
That ratio is worth keeping in view.
\end{SyntaxBox}

\begin{ExplainBox}{A network's result is a random variable}
Weight initialization and batch order are random, so \c{random_state} changes the
fitted model. Refitting this network with seeds 1 through 8 gives test ROC-AUC
between 0.895 and 0.916, mean 0.905. That spread is comparable to the difference
between the models being compared, so a single run is not evidence that one
architecture beat another. Report the seed and, when a comparison is close,
several seeds -- the same discipline \booklink{ch:20}{Chapter 20} demanded of
metaheuristics.
\end{ExplainBox}

\hh{Width and regularization control capacity}

\bookfigure[0.98\textwidth]{Figures/Generated/ch22_network_capacity.pdf}
{Training ROC-AUC climbs from 0.763 to 0.996 as the hidden layer widens while cross-validated ROC-AUC peaks near 16 neurons and then declines; the learning curve still rises at 540 rows, so more data would help more than more neurons.}
{fig:ch22-capacity}

\begin{center}
\begin{tabular}{llll}
\textbf{Hidden neurons} & \textbf{Training} & \textbf{Validation} & \textbf{Validation s.d.} \\
2 & 0.763 & 0.741 & 0.072 \\
4 & 0.872 & 0.837 & 0.043 \\
16 & 0.905 & 0.837 & 0.042 \\
64 & 0.973 & 0.820 & 0.031 \\
128 & 0.996 & 0.802 & 0.048 \\
\end{tabular}
\end{center}

This is the same shape as the tree-depth curve in
\booklink{ch:21}{Chapter 21}, and it carries the same lesson: the training score
rises with capacity forever, and only the validation score identifies the useful
setting. Two neurons underfit; 128 neurons reach 0.996 on the training folds and
0.802 on held-out data.

The curve is not perfectly smooth -- 8 neurons validates slightly below 4 and 16.
With a fold-to-fold standard deviation near 0.04, that dip is inside the noise.
Each width is also a separate optimization run with its own initialization, so
some of the wobble is the optimizer rather than the architecture. Do not read a
0.03 difference as a finding when the spread is 0.04.

\hh{Learning curves say whether more data would help}

A validation curve varies the model. A learning curve varies the amount of data.

\begin{center}
\begin{tabular}{lll}
\textbf{Training rows used} & \textbf{Training ROC-AUC} & \textbf{Validation ROC-AUC} \\
108 & 1.000 & 0.736 \\
216 & 0.939 & 0.807 \\
324 & 0.953 & 0.823 \\
432 & 0.926 & 0.841 \\
540 & 0.905 & 0.837 \\
\end{tabular}
\end{center}

With 108 rows the network memorizes the training data perfectly and generalizes
poorly: a gap of 0.26. As rows are added the training score falls, the validation
score rises, and the gap narrows to about 0.07. The validation curve has not yet
flattened, which is the diagnostic: on this problem, collecting more rows would
buy more than adding neurons.

\hh{Compare against the earlier models honestly}

\begin{center}
\begin{tabular}{lll}
\textbf{Model} & \textbf{Cross-validated ROC-AUC} & \textbf{Test ROC-AUC} \\
Logistic pipeline & 0.766 +/- 0.071 & 0.776 \\
Random forest & 0.836 +/- 0.041 & 0.911 \\
Neural network & 0.837 +/- 0.042 & 0.895 \\
\end{tabular}
\end{center}

The network decisively beats logistic regression, because it can represent the
two-sided setpoint effect that a straight line cannot. Against the random forest
it is a tie: 0.837 against 0.836 in cross-validation, with fold spreads of about
0.04 that overlap almost completely.

Watch what that tie does. Selecting purely by cross-validated score picks the
network, by one thousandth of a point -- and on the protected test set the forest
then scores higher, 0.911 against 0.895. Neither number proves the forest is
better either. Both observations are the same lesson: a gap of 0.001 against a
spread of 0.04 is not a result, and an automatic \c{max()} over candidate scores
will happily report it as one.

That tie is the honest headline of this chapter. The forest reached its score
with three hyperparameters and no scaling requirement; the network needed
standardization, a width, a regularization strength, an iteration cap, and a
seed, and it repays none of that effort here. On the test set the network does
find more complaints -- 34 of 54 against the forest's 31, because its probability
distribution sits differently relative to the 0.50 threshold -- but with 54
positive rows that difference is well inside sampling noise.

\begin{ExplainBox}{Neural networks earn their keep on structure, not on tables}
Networks dominate where the input has structure a human would otherwise have to
hand-engineer: neighbouring pixels in an image, ordered samples in audio, tokens
in text. There, learned intermediate representations replace features nobody
knows how to write down. On a modest table of six meaningful engineered columns,
gradient-boosted trees and random forests remain extremely strong, and they are
cheaper to tune, more stable across seeds, and easier to explain. Choosing a
network because it is the most advanced tool available is the same error as
choosing a metaheuristic for a smooth convex objective.
\end{ExplainBox}

\hh{When a neural network is the wrong choice}

\begin{itemize}
    \item \textbf{Few rows.} Hundreds of samples rarely support thousands of
          parameters.
    \item \textbf{Tabular data with meaningful columns.} Tree ensembles are a
          strong default and usually the right first attempt.
    \item \textbf{A hard interpretability requirement.} If a decision must be
          explained to an occupant, a regulator, or a review board, start with a
          model whose logic can be stated.
    \item \textbf{No tuning budget.} A network has more hyperparameters that
          matter, and its untuned performance is often poor.
    \item \textbf{A linear relationship.} If the effect really is a weighted sum,
          the extra capacity only adds variance.
\end{itemize}

\hh{Applied example: a network in the Chapter 21 workflow}
\label{sec:c22-application}

The runnable program \c{examples/c22_neural_network_comfort.py} trains the
hand-written network with a verified gradient, then adds \c{MLPClassifier} as a
fourth candidate in the Chapter 21 comparison and exports the loss curve and
decision-boundary figure.

\begin{lstlisting}[
    language=python,
    style=block,
    caption={The network takes its turn as one candidate among several},
    label={c22_application},
    captionpos=b
]
from sklearn.ensemble import RandomForestClassifier
from sklearn.linear_model import LogisticRegression
from sklearn.model_selection import StratifiedKFold, cross_val_score

folds = StratifiedKFold(n_splits=5, shuffle=True, random_state=42)

candidates = {
    "logistic pipeline": Pipeline([
        ("scale", StandardScaler()),
        ("model", LogisticRegression(max_iter=1000)),
    ]),
    "random forest": RandomForestClassifier(n_estimators=300, random_state=42),
    "neural network": network,
}

for name, estimator in candidates.items():
    scores = cross_val_score(estimator, X_train, y_train, cv=folds, scoring="roc_auc")
    print(f"{name:<20} ROC-AUC {scores.mean():.3f} +/- {scores.std():.3f}")

# Only after the comparison is settled does the test set open, once.
best_name, best_estimator = max(
    candidates.items(),
    key=lambda item: cross_val_score(
        item[1], X_train, y_train, cv=folds, scoring="roc_auc"
    ).mean(),
)
print("Selected by cross-validation:", best_name)
\end{lstlisting}

\begin{SyntaxBox}{Nothing about the workflow changed}
The network is a \c{Pipeline} with \c{fit()}, \c{predict()}, and
\c{predict_proba()}, so it slots into the Chapter 21 procedure without special
handling. The split, the folds, the baseline, the scoring metric, and the
single test evaluation are all unchanged. A more powerful model family does not
earn relaxed evidence standards; if anything, its extra tuning freedom makes
the protected test set more important.
\end{SyntaxBox}

\hh{Common mistakes and useful diagnoses}

\begin{itemize}
    \item \textbf{Backpropagation is never gradient-checked.} A wrong gradient
          still trains, badly. Compare against finite differences.
    \item \textbf{Features are not scaled.} Cross-validated ROC-AUC fell from
          0.837 to 0.617 in this chapter's data.
    \item \textbf{All weights initialized to zero.} Every hidden neuron then
          computes the same function and receives the same gradient forever.
    \item \textbf{The loss becomes \c{nan}.} Usually \c{log(0)}; clip
          probabilities away from 0 and 1, and check for too large a learning
          rate.
    \item \textbf{A \c{ConvergenceWarning} is suppressed.} It means training
          stopped before it finished; raise \c{max_iter} or fix the scaling.
    \item \textbf{\c{random_state} is unset.} Two runs then disagree and neither
          is reproducible.
    \item \textbf{Architecture is chosen on the test set.} Width and \c{alpha}
          are hyperparameters; select them with cross-validation.
    \item \textbf{A small difference is read as an improvement.} Compare the gap
          against the fold spread and the seed-to-seed spread.
    \item \textbf{The network is assumed to be best because it is a network.}
          Benchmark against a tree ensemble and a linear model every time.
\end{itemize}

\hh{Practice ladder}

\hhh{Level 0 -- Identify}

In Listing~\ref{c22_forward}, identify the shape of every array for 675 training
rows, 6 features, and 8 hidden neurons, and count the parameters.

\hhh{Level 1 -- Modify}

Retrain the hand-written network with \c{hidden_count=4} and \c{hidden_count=16}
at the same learning rate. Report the final training loss and test ROC-AUC for
each, and say whether the differences exceed what you would expect from changing
the seed.

\hhh{Level 2 -- Explain}

The hidden layer uses tanh but the output uses sigmoid. Explain why the output
cannot use tanh for this task, and why the hidden layer would train poorly if it
used sigmoid.

\hhh{Level 3 -- Build}

Write \c{train_with_validation(X, targets, X_valid, targets_valid, patience)}
that records both losses each epoch and stops when the validation loss has failed
to improve for \c{patience} consecutive epochs. Return the parameters from the
best epoch, not the last one.

\hhh{Level 4 -- Challenge}

Add a second hidden layer to the hand-written network: extend \c{initialize},
\c{forward}, and \c{backward}, and gradient-check the result before training.
Then compare the two-layer network against the one-layer network across five
seeds and state whether the extra layer is justified on this dataset.

\hh{Practice solutions}

\textbf{Level 0:} \c{X} is (675, 6); \c{W1} is (6, 8) and \c{b1} is (8,);
\c{hidden_output} is (675, 8); \c{W2} is (8, 1) and \c{b2} is (1,); the returned
probabilities are (675,). The parameter count is
$6 \times 8 + 8 + 8 \times 1 + 1 = 65$.

\textbf{Level 2:} The output must be a probability in $(0, 1)$, and cross-entropy
takes the logarithm of that output and of one minus it. Tanh returns values in
$(-1, 1)$, so negative outputs would make both logarithms undefined. In the
hidden layer, sigmoid's slope never exceeds 0.25, so gradients shrink by at least
a factor of four per layer on the way back; tanh is zero-centred with a maximum
slope of 1 and passes signal back far better.

\textbf{Level 3}

\begin{lstlisting}[
    language=python,
    style=block,
    caption={Training with early stopping on a held-out loss},
    label={c22_solution_early_stopping},
    captionpos=b
]
import numpy as np


def train_with_validation(X, targets, X_valid, targets_valid, patience=20,
                          hidden_count=8, learning_rate=0.5, epochs=2000, seed=0):
    parameters = initialize(X.shape[1], hidden_count, seed)
    best_parameters = {key: value.copy() for key, value in parameters.items()}
    best_loss = np.inf
    waited = 0
    history = []

    for _ in range(epochs):
        hidden_output, probabilities = forward(parameters, X)
        gradients = backward(parameters, X, targets, hidden_output, probabilities)
        for key in parameters:
            parameters[key] -= learning_rate * gradients[key]

        validation_loss = cross_entropy(forward(parameters, X_valid)[1], targets_valid)
        history.append((cross_entropy(probabilities, targets), validation_loss))

        if validation_loss < best_loss - 1e-6:
            best_loss = validation_loss
            best_parameters = {key: value.copy() for key, value in parameters.items()}
            waited = 0
        else:
            waited += 1
            if waited >= patience:
                break

    return best_parameters, np.array(history), best_loss
\end{lstlisting}

The parameters are copied when the best validation loss is seen, because training
continues past that point and would otherwise overwrite them. The held-out set
used for stopping must not be the protected test set; take it from the training
data.

\hh{Chapter checkpoint}

\begin{enumerate}
    \item What two operations make up one neuron?
    \item Why does a stack of linear layers with no activation gain nothing?
    \item What does the slope of an activation function control during training?
    \item Why is cross-entropy the training loss rather than accuracy?
    \item What does backpropagation compute, and in which direction?
    \item How do you verify a backpropagation implementation?
    \item Why must all weights not be initialized to zero?
    \item What does feature scaling do for a network, quantitatively, on this data?
    \item How do a validation curve and a learning curve differ in what they vary?
    \item Name two situations in which a tree ensemble is the better choice.
\end{enumerate}

\hh{Checkpoint answers}

\begin{enumerate}
    \item A weighted sum of its inputs plus a bias, then a nonlinear activation.
    \item The weight matrices multiply into a single matrix, so the composition
          expresses exactly what one linear layer expresses.
    \item How much gradient signal reaches earlier layers; a slope near zero
          stalls learning.
    \item Cross-entropy is differentiable and penalizes confident errors, while
          accuracy is a flat step function that provides no gradient.
    \item The derivative of the loss with respect to every parameter, computed
          from the output layer backward using the chain rule.
    \item Compare each analytic derivative against a central finite-difference
          estimate; the largest gap should be near \c{1e-10}.
    \item Identical weights make every hidden neuron compute the same function
          and receive the same gradient, so they never differentiate.
    \item It raised cross-validated ROC-AUC from 0.617 to 0.837 and test ROC-AUC
          from 0.681 to 0.895.
    \item A validation curve varies a model hyperparameter; a learning curve
          varies the number of training rows.
    \item Modest tabular data with meaningful columns, and any situation
          requiring interpretability or a small tuning budget.
\end{enumerate}

\begin{tcolorbox}[title=Part VII summary,colframe=orange,colback=white,arc=2mm]
Algorithms turn specifications into finite methods whose correctness and growth
can be explained. Optimization searches a modeled feasible space: gradient
methods descend into the nearest valley, metaheuristics trade evaluations for
the ability to escape it, and both require an equal-budget baseline and an
honest account of randomness. Machine learning estimates rules from data, and
model selection is itself a search that must run inside training data so the
test set can still answer the only question that matters. A neural network is
those two threads tied together -- a layered model fitted by gradient descent and
selected by cross-validation -- and it is subject to exactly the same evidence
standards as everything before it. Across all six chapters the habit is the
same: define the contract, establish a baseline, trace the computation, verify
the gradient or the result, measure with more than one number, report the
spread, and state plainly what the result does not prove.
\end{tcolorbox}
